%% kapa: exam=1 round=25 year=2014 period=1 form=A subject=economics subject_ko=경제원론 seq=2 number=42
%% source: Q-Net 감정평가사 기출문제 (cst00309) artlSeq=5101207 file=extracted/1008/5101207/2014년 제25회 감정평가사 1차시험 기출문제/1교시(A형).hwp
%% conversion: kordoc 4.24.51 -> tools/gen1.py (auto); review=proofread-sonnet; flags=tikz
\begin{question}{42}
다음 그림은 $X$재에 대한 수요곡선이다. 다음 설명 중 옳은 것을 모두 고른 것은? (단, $X$재는 정상재이다.)
\begin{center}
\begin{tikzpicture}[font=\footnotesize]
  \draw[->] (0,-0.2) -- (0,3.4) node[left] {P};
  \draw[->] (-0.2,0) -- (4.2,0) node[right] {Q};
  \draw (0,3) -- (3,0);
  \node[left] at (0,3) {100};
  \node[below] at (3,0) {100};
\end{tikzpicture}
\end{center}
\begin{framedtext}
ㄱ. 가격이 30원이면 $X$재의 수요량은 70이다. \par
ㄴ. 가격에 상관없이 가격탄력성의 크기는 일정하다. \par
ㄷ. $X$재의 시장이 독점시장이라면 독점기업이 이윤극대화를 할 때 설정하는 가격은 50원 이상이다. \par
ㄹ. 소득이 증가하는 경우 수요곡선은 좌측으로 이동한다. \par
\end{framedtext}
\begin{choices}
  \item ㄱ, ㄴ
  \item ㄱ, ㄷ
  \item ㄴ, ㄷ
  \item ㄴ, ㄹ
  \item ㄷ, ㄹ
\end{choices}
\end{question}
\answer{2}
